\documentclass[12pt,a4paper]{article}

\usepackage[utf8]{inputenc}
\usepackage[brazil]{babel}
\usepackage[T1]{fontenc}
\usepackage{geometry}
\usepackage{setspace}
\usepackage{indentfirst}
\usepackage{listings}
\usepackage{xcolor}
\usepackage{tikz}
\usetikzlibrary{shapes.geometric, arrows.meta, positioning}

\geometry{margin=2.5cm}
\onehalfspacing

% Definição de cores acadêmicas
\definecolor{ludiblue}{RGB}{30, 100, 180}
\definecolor{ludipurple}{RGB}{120, 50, 160}
\definecolor{ludigray}{RGB}{245, 245, 245}
\definecolor{ludigreen}{RGB}{0, 120, 60}
\definecolor{ludicomment}{RGB}{100, 100, 100}

% Configuração do listings para a linguagem LudiCode real
\lstdefinelanguage{LudiCode}{
  keywords={var, se, senao, enquanto, repita, funcao, mostrar, verdadeiro, falso, retorne},
  keywordstyle=\color{ludiblue}\bfseries,
  ndkeywords={},
  sensitive=true,
  comment=[l]{//},
  morecomment=[l]{\#},
  commentstyle=\color{ludicomment}\itshape,
  stringstyle=\color{ludigreen},
  morestring=[b]",
}

\lstset{
    language=LudiCode,
    basicstyle=\ttfamily\small,
    backgroundcolor=\color{ludigray},
    breaklines=true,
    frame=single,
    rulecolor=\color{ludipurple!40},
    numbers=left,
    numberstyle=\tiny\color{gray},
    keywordstyle=\color{ludiblue}\bfseries,
    stringstyle=\color{ludigreen},
    commentstyle=\color{ludicomment}\itshape,
    showstringspaces=false,
    tabsize=4,
    captionpos=b,
}

% Estilo TikZ para o diagrama de fluxo
\tikzset{
    bloco/.style={
        rectangle, rounded corners=4pt,
        minimum width=2.2cm, minimum height=0.9cm,
        text centered, font=\small\bfseries,
        draw=ludiblue, fill=ludiblue!12, text=ludiblue
    },
    seta/.style={-Stealth, thick, color=ludipurple}
}

\begin{document}

% ────────────────────────────────────────────────
%  CAPA
% ────────────────────────────────────────────────
\begin{titlepage}
\begin{center}

\textbf{CENTRO UNIVERSITÁRIO MAURÍCIO DE NASSAU}

\vspace{2.5cm}

\textbf{\Large DESENVOLVIMENTO DO COMPILADOR LUDICODE}

\vspace{0.5cm}
\textbf{\large Uma Linguagem de Programação Estruturada em Português}

\vspace{2cm}

\textbf{Equipe de Desenvolvimento:}

\vspace{0.5cm}

Abraão Araujo dos Santos\\
Erick Alves de Souza\\
Evandro França da Silva Filho\\
Gabriel Marques Barbosa de Santana\\
Gabriella Maria Nascimento da Silva\\
Jonnas Kauan Pereira Novaes\\
Mateus De Miranda Santos Moura\\
Matheus Pinheiro da Cunha\\

\vfill

Maceió -- AL\\
2026

\end{center}
\end{titlepage}

\tableofcontents
\newpage

% ────────────────────────────────────────────────
%  1. INTRODUÇÃO
% ────────────────────────────────────────────────
\section{Introdução}

Para que o computador compreenda as instruções, o projeto utiliza uma arquitetura clássica de compiladores, desenvolvida em Python com o auxílio da biblioteca SLY. O processo funciona como uma linha de montagem: primeiro, o sistema lê o texto e o organiza em pedaços lógicos (análise léxica); em seguida, ele verifica se a estrutura dessas frases obedece às regras da linguagem para montar um mapa de execução, chamado de Árvore Sintática Abstrata (análise sintática); por fim, ele valida se o comando faz sentido matematicamente e logicamente (validação semântica) antes de efetivamente executá-lo (interpretador).

% ────────────────────────────────────────────────
%  2. ESTRUTURA E RECURSOS DA LINGUAGEM
% ────────────────────────────────────────────────

\section{Estrutura e Recursos da Linguagem}

A linguagem LudiCode possui sintaxe imperativa e estruturada por blocos delimitados por chaves (\texttt{\{\}}). No estado atual do código-fonte, a parte principal da linguagem está implementada sobre declarações, atribuições, expressões, comandos de saída e estruturas de controle. O interpretador também possui uma camada de simulação para recursos de robótica educacional, mantendo estado de sensores, motores e eventos, embora esses comandos de hardware ainda não estejam ligados diretamente a uma produção sintática geral no arquivo \texttt{parse.py}.

Os principais recursos efetivamente presentes na versão atual são:

\begin{itemize}
    \item \textbf{Declaração e atribuição dinâmica:} uso da palavra reservada \texttt{var} para criação de variáveis na tabela de símbolos e uso de \texttt{=} para atualização de variáveis já declaradas.
    \item \textbf{Tipos básicos:} suporte a números inteiros, números decimais, textos entre aspas e valores booleanos \texttt{verdadeiro} e \texttt{falso}.
    \item \textbf{Expressões aritméticas e relacionais:} suporte aos operadores \texttt{+}, \texttt{-}, \texttt{*}, \texttt{/}, \texttt{==}, \texttt{<}, \texttt{>}, \texttt{<=} e \texttt{>=} dentro das regras do parser.
    \item \textbf{Controle de fluxo condicional:} estruturas \texttt{se} e \texttt{senao}, com condição entre parênteses e corpo estruturado em bloco.
    \item \textbf{Laços de repetição:} suporte ao laço condicional \texttt{enquanto} e ao comando \texttt{repita} seguido de um número literal e um bloco.
    \item \textbf{Saída em terminal:} comando nativo \texttt{mostrar()} implementado no parser e no interpretador para exibição de valores.
    \item \textbf{Simulação de robótica no interpretador:} a classe \texttt{SimulationRuntime} mantém motores, sensores e eventos; o interpretador possui funções internas para \texttt{ligarMotor}, \texttt{desligarMotor}, \texttt{lerSensor} e \texttt{configurarSensor}, acessíveis quando a AST normalizada contém nós de chamada.
\end{itemize}

O exemplo de listagem a seguir ilustra o comportamento típico de validação condicional em código-fonte LudiCode:

\begin{lstlisting}[caption={Exemplo de estrutura de controle condicional}]
var idade = 18

se (idade >= 18) {
    mostrar("Maior de idade")
} senao {
    mostrar("Menor de idade")
}
\end{lstlisting}

% ────────────────────────────────────────────────
%  3. GRAMÁTICA FORMAL (BNF)
% ────────────────────────────────────────────────

\section{Gramática Formal da Linguagem}

A sintaxe da linguagem LudiCode é especificada formalmente por meio da Notação Backus-Naur (\textbf{BNF}), descrevendo as produções aceitas pelo analisador sintático implementado em \texttt{parse.py}. A gramática abaixo documenta o comportamento atualmente codificado no parser.

\begin{lstlisting}[language={}, caption={Especificação BNF da Linguagem LudiCode conforme parse.py}, keywordstyle=\color{black}, numbers=none, frame=single]
<program>      ::= <statement>
               | <program> <statement>

<statement>    ::= "var" <id> "=" <expr>
               | <id> "=" <expr>
               | "mostrar" "(" <expr> ")"
               | "se" "(" <expr> ")" <bloco>
               | "se" "(" <expr> ")" <bloco> "senao" <bloco>
               | "se" "(" <expr> ")" <bloco> "senao" <statement>
               | "se" "(" <expr> ")" <statement> "senao" <bloco>
               | "enquanto" "(" <expr> ")" <bloco>
               | "repita" <numero> <bloco>

<bloco>        ::= "{" <program> "}"
               | "{" <statement> "}"

<expr>         ::= <expr> "+" <expr>
               | <expr> "-" <expr>
               | <expr> "*" <expr>
               | <expr> "/" <expr>
               | <expr> "==" <expr>
               | <expr> "<"  <expr>
               | <expr> ">"  <expr>
               | <expr> "<=" <expr>
               | <expr> ">=" <expr>
               | <numero>
               | <texto>
               | "verdadeiro"
               | "falso"
               | <id>

<numero>       ::= [0-9]+ | [0-9]+ "." [0-9]+
<texto>        ::= '"' <qualquer-caractere>* '"'
<id>           ::= [a-zA-Z_][a-zA-Z0-9_]*
\end{lstlisting}

É importante observar que os tokens \texttt{FUNCAO} e \texttt{RETORNE} já estão reconhecidos no analisador léxico, porém não possuem produções sintáticas associadas no \texttt{parse.py} desta versão. Da mesma forma, as chamadas genéricas de função usadas internamente pelo interpretador para recursos de robótica ainda não aparecem como regra gramatical geral no parser.

% ────────────────────────────────────────────────
%  4. DESENVOLVIMENTO DO SISTEMA
% ────────────────────────────────────────────────
\section{Desenvolvimento do Sistema}

A implementação foi dividida modularmente em quatro arquivos funcionais independentes, cujas atribuições respeitam o pipeline tradicional de compilação.

\subsection{Pipeline de Processamento}

O fluxo de processamento de dados ocorre de maneira síncrona e hierárquica, transformando a sequência de caracteres textuais brutos em uma execução semântica controlada, conforme ilustrado no diagrama a seguir:

\begin{center}
\begin{tikzpicture}[node distance=0.6cm and 1.2cm]

  \node[bloco] (src)  {\shortstack{Código-Fonte\\.ludi}};
  \node[bloco, right=of src]  (lex)  {\shortstack{Lexer\\(Tokens)}};
  \node[bloco, right=of lex]  (par)  {Parser};
  \node[bloco, right=of par]  (ast)  {AST};
  \node[bloco, right=of ast]  (int)  {Interpretador};

  \draw[seta] (src) -- (lex);
  \draw[seta] (lex) -- node[above,font=\tiny]{tokens} (par);
  \draw[seta] (par) -- node[above,font=\tiny]{árvore} (ast);
  \draw[seta] (ast) -- (int);

\end{tikzpicture}
\end{center}


\subsection{Análise Léxica (\texttt{lexer.py})}

O analisador léxico (\texttt{LudiCodeLexer}) converte o fluxo de caracteres de entrada em uma sequência de \textit{tokens}. As palavras reservadas são reconhecidas por meio da regra geral de identificadores e, em seguida, remapeadas para tokens específicos por um dicionário interno.

A tabela a seguir consolida os principais tokens implementados no arquivo \texttt{lexer.py}:

\begin{center}
\begin{tabular}{llp{7cm}}
\hline
\textbf{Token Léxico} & \textbf{Padrão/Palavra} & \textbf{Descrição Funcional} \\
\hline
\texttt{VAR}        & \texttt{var}       & Declaração de variável. \\
\texttt{SE}         & \texttt{se}        & Início de estrutura condicional. \\
\texttt{SENAO}      & \texttt{senao}     & Bloco alternativo da estrutura condicional. \\
\texttt{ENQUANTO}   & \texttt{enquanto}  & Laço de repetição com condição inicial. \\
\texttt{REPITA}     & \texttt{repita}    & Laço de repetição com quantidade numérica literal. \\
\texttt{FUNCAO}     & \texttt{funcao}    & Palavra reservada reconhecida pelo lexer, ainda sem regra sintática no parser. \\
\texttt{MOSTRAR}    & \texttt{mostrar}   & Comando nativo de impressão em console. \\
\texttt{VERDADEIRO} & \texttt{verdadeiro}& Literal lógico booleano positivo. \\
\texttt{FALSO}      & \texttt{falso}     & Literal lógico booleano negativo. \\
\texttt{RETORNE}    & \texttt{retorne}   & Palavra reservada reconhecida pelo lexer, ainda sem regra sintática no parser. \\
\texttt{ID}         & Expressão regular  & Identificadores genéricos de variáveis. \\
\texttt{NUMERO}     & Inteiro ou decimal & Valores numéricos convertidos para \texttt{int} ou \texttt{float}. \\
\texttt{TEXTO}      & Texto entre aspas  & Cadeias textuais, com remoção das aspas na etapa léxica. \\
\texttt{IGUAL}      & \texttt{==}        & Operador relacional de igualdade. \\
\texttt{ATRIBUICAO} & \texttt{=}         & Operador de atribuição. \\
\texttt{MENOR}      & \texttt{<}         & Operador relacional menor que. \\
\texttt{MAIOR}      & \texttt{>}         & Operador relacional maior que. \\
\texttt{MENOR\_IGUAL} & \texttt{<=}      & Operador relacional menor ou igual. \\
\texttt{MAIOR\_IGUAL} & \texttt{>=}      & Operador relacional maior ou igual. \\
\hline
\end{tabular}
\end{center}

Além dos tokens nomeados, o lexer define como literais os símbolos \texttt{+}, \texttt{-}, \texttt{*}, \texttt{/}, \texttt{(}, \texttt{)}, \texttt{\{}, \texttt{\}}, \texttt{,}, \texttt{.} e \texttt{;}. Comentários de linha iniciados por \texttt{//} ou \texttt{\#} são ignorados. Quebras de linha atualizam o contador interno \texttt{lineno}, utilizado nas mensagens de erro léxico.


\subsection{Análise Sintática (\texttt{parse.py}) e Árvore Abstrata (AST)}

A análise sintática baseia-se em regras de produção LALR(1) parametrizadas por decoradores da biblioteca SLY. A resolução de ambiguidades gramaticais de expressões é estabelecida pela tupla de precedência presente no próprio parser:

\begin{lstlisting}[language=Python, caption={Declaração de precedência estática no parser}]
precedence = (
    ('nonassoc', 'IGUAL', 'MENOR', 'MAIOR', 'MENOR_IGUAL', 'MAIOR_IGUAL'),
    ('left', '+', '-'),
    ('left', '*', '/'),
)
\end{lstlisting}

Ao processar a entrada, o parser sintetiza uma Árvore Sintática Abstrata (\textbf{AST}) estruturada em \textbf{tuplas aninhadas} nativas do Python. Essa representação é compacta e compatível com o método \texttt{normalize()} do interpretador, que também aceita, em alguns casos, nós no formato de dicionário com a chave \texttt{tipo}.

Abaixo estão listados os formatos estruturais das tuplas geradas pelo parser para as principais construções sintáticas:

\begin{center}
\begin{tabular}{lp{8.5cm}}
\hline
\textbf{Construção Sintática} & \textbf{Estrutura do Nó na AST (Tupla)} \\
\hline
Declaração de variável & \texttt{('declaracao\_var', nome, expressao)} \\
Atribuição simples     & \texttt{('atribuicao', nome, expressao)} \\
Comando Mostrar        & \texttt{('mostrar', expressao)} \\
Condicional simples    & \texttt{('se', condicao, bloco\_verdadeiro)} \\
Condicional composta   & \texttt{('se', condicao, bloco\_verdadeiro, bloco\_falso)} \\
Laço Enquanto          & \texttt{('enquanto', condicao, bloco\_corpo)} \\
Laço Repita            & \texttt{('repita', numero, bloco\_corpo)} \\
Operação binária       & \texttt{('operacao', operador, esq, dir)} \\
Número                 & \texttt{('numero', valor)} \\
Texto                  & \texttt{('texto', valor)} \\
Booleano               & \texttt{('booleano', True/False)} \\
Variável               & \texttt{('variavel', nome)} \\
\hline
\end{tabular}
\end{center}

No estado atual, o parser implementa diretamente o comando \texttt{mostrar()}, mas não possui uma regra geral para chamadas como \texttt{nomeDaFuncao(...)}. Por isso, os recursos de hardware presentes no interpretador aparecem como suporte interno de execução, e não como comandos completamente descritos pela gramática textual desta versão.


\subsection{Interpretador Executável (\texttt{interpreter.py})}

O interpretador (\texttt{LudiCodeInterpreter}) percorre a AST normalizada e executa cada nó de acordo com uma tabela interna de executores. A tabela de símbolos é implementada pela classe \texttt{SymbolTable}, responsável por armazenar variáveis declaradas, impedir dupla declaração e bloquear leitura ou atribuição de identificadores inexistentes.

A execução é organizada por três componentes principais:

\begin{itemize}
    \item \textbf{\texttt{SymbolTable}:} mantém o dicionário \texttt{values}, com as variáveis e seus valores atuais.
    \item \textbf{\texttt{SimulationRuntime}:} mantém o estado de \texttt{sensors}, \texttt{motors} e \texttt{events}, servindo como camada de simulação para robótica educacional.
    \item \textbf{\texttt{LudiCodeInterpreter}:} normaliza a AST, avalia expressões, executa comandos, controla laços e retorna o estado final da execução.
\end{itemize}

O método \texttt{run()} retorna um dicionário contendo quatro grupos de informação: \texttt{saida}, \texttt{variaveis}, \texttt{sensores}, \texttt{motores} e \texttt{eventos}. Esse retorno é utilizado pelo arquivo \texttt{main.py} para imprimir o resultado final da execução de arquivos \texttt{.ludi}.

O interpretador possui suporte aos operadores \texttt{+}, \texttt{-}, \texttt{*}, \texttt{/}, \texttt{==}, \texttt{!=}, \texttt{<}, \texttt{>}, \texttt{<=}, \texttt{>=}, \texttt{e} e \texttt{ou}. Entretanto, no parser atual, apenas os operadores aritméticos e relacionais definidos em \texttt{parse.py} são produzidos diretamente a partir do código-fonte textual.

Também há funções internas registradas no dicionário \texttt{builtins}: \texttt{ligarmotor}, \texttt{desligarmotor}, \texttt{lersensor}, \texttt{configurarsensor} e \texttt{mostrar}. Essas funções executam ações sobre a camada \texttt{SimulationRuntime} ou sobre a saída do programa. No caso de \texttt{mostrar}, além de armazenar o valor em \texttt{output}, o método também imprime o valor imediatamente no terminal.

O laço \texttt{enquanto} é protegido por um limite de segurança configurado pelo atributo \texttt{loop\_limit}, cujo valor padrão atual é \texttt{1000}. Caso esse limite seja ultrapassado, o interpretador lança uma exceção semântica para evitar execução infinita.

% ────────────────────────────────────────────────
%  5. TRATAMENTO DE ERROS E ANÁLISE SEMÂNTICA
% ────────────────────────────────────────────────

\section{Tratamento de Erros e Análise Semântica}

O sistema valida a semântica do programa durante a interpretação dos nós da AST. Quando alguma regra é violada, o interpretador lança uma exceção customizada do tipo \texttt{SemanticError}, com mensagem direta para facilitar a compreensão do problema.

As principais validações semânticas presentes no código atual são:

\begin{enumerate}
    \item \textbf{Dupla declaração de variável:} ocorre quando \texttt{var} tenta declarar um identificador já existente na \texttt{SymbolTable}.
    \item \textbf{Variável não declarada:} ocorre ao tentar ler ou atribuir valor a uma variável que ainda não foi criada com \texttt{var}.
    \item \textbf{Divisão por zero:} interceptada durante a avaliação do operador \texttt{/}, convertendo o erro nativo do Python em \texttt{SemanticError}.
    \item \textbf{Incompatibilidade de tipos:} ocorre quando uma operação recebe valores incompatíveis, por exemplo, tentar subtrair um número de um texto.
    \item \textbf{Operador não suportado:} ocorre quando a AST contém um operador não registrado no dicionário \texttt{operators}.
    \item \textbf{Função não suportada:} ocorre quando uma chamada de função não aparece no dicionário \texttt{builtins}.
    \item \textbf{Sensor não configurado:} ocorre quando \texttt{lerSensor} tenta consultar um sensor inexistente no estado de simulação.
    \item \textbf{Limite de laço excedido:} ocorre quando o laço \texttt{enquanto} ultrapassa o valor definido em \texttt{loop\_limit}.
\end{enumerate}

No arquivo \texttt{main.py}, erros léxicos capturados no modo interativo são tratados com mensagem específica. Na execução por arquivo, erros sintáticos interrompem a execução quando o parser retorna uma AST nula, enquanto erros de execução são exibidos na forma \texttt{Erro durante a execucao: ...}.

% ────────────────────────────────────────────────
%  6. TESTES E LOGS REALIZADOS
% ────────────────────────────────────────────────

\section{Testes Executados e Saídas do Sistema}

Os scripts presentes na pasta \texttt{tests/} validam estruturas de repetição, desvio condicional e tratamento de erros semânticos. A execução ocorre pelo arquivo \texttt{main.py}, utilizando o comando \texttt{python main.py caminho/do/arquivo.ludi}. Durante a geração do parser, a biblioteca SLY pode emitir avisos técnicos sobre tokens definidos e ainda não utilizados, como \texttt{FUNCAO} e \texttt{RETORNE}; esses avisos não impedem a execução dos testes listados abaixo.

\subsection{Teste do Laço \texttt{enquanto}}

O script \texttt{teste\_enquanto.ludi} avalia uma variável contadora incremental executada três vezes consecutivas.

\begin{lstlisting}[caption={tests/teste\_enquanto.ludi}]
var x = 0

enquanto (x < 3) {
    mostrar(x)
    x = x + 1
}
\end{lstlisting}

\textbf{Saída funcional gerada no terminal:}
\begin{lstlisting}[language={}, numbers=none, frame=single, caption={Saída do console para o laço enquanto}]
0
1
2

Saida do programa:
0
1
2

Estado final:
Variaveis: {'x': 3}
Sensores: {}
Motores: {}
Eventos: []
\end{lstlisting}

Nesse teste, os valores \texttt{0}, \texttt{1} e \texttt{2} aparecem duas vezes porque o método \texttt{mostrar} imprime imediatamente no terminal e, ao final da execução, o \texttt{main.py} percorre novamente a lista \texttt{saida} para exibir o resumo do programa.

\subsection{Teste Condicional Estruturado}

O arquivo \texttt{teste\_se.ludi} valida uma estrutura \texttt{se/senao} com comparação maior ou igual.

\begin{lstlisting}[caption={tests/teste\_se.ludi}]
var idade = 18

se (idade >= 18) {
    mostrar("Maior de idade")
} senao {
    mostrar("Menor de idade")
}
\end{lstlisting}

\textbf{Saída funcional gerada no terminal:}
\begin{lstlisting}[language={}, numbers=none, frame=single, caption={Saída do console para desvio condicional}]
Maior de idade

Saida do programa:
Maior de idade

Estado final:
Variaveis: {'idade': 18}
Sensores: {}
Motores: {}
Eventos: []
\end{lstlisting}

\subsection{Exemplos de Captura de Erros Semânticos}

O arquivo \texttt{teste\_tipos.ludi} valida o tratamento de incompatibilidade de tipos durante uma operação aritmética.

\begin{lstlisting}[caption={tests/teste\_tipos.ludi}]
var x = "abc" - 1
\end{lstlisting}
\begin{lstlisting}[language={}, numbers=none, frame=single, caption={Log de erro para tipos incompatíveis}]
Erro durante a execucao: Tipos incompativeis para o operador '-': str e int.
\end{lstlisting}

O arquivo \texttt{teste\_divisao\_zero.ludi} valida o tratamento de divisão por zero.

\begin{lstlisting}[caption={tests/teste\_divisao\_zero.ludi}]
var x = 10 / 0
\end{lstlisting}
\begin{lstlisting}[language={}, numbers=none, frame=single, caption={Log de erro para divisão por zero}]
Erro durante a execucao: Divisao por zero.
\end{lstlisting}

% ────────────────────────────────────────────────
%  7. CONCLUSÃO
% ────────────────────────────────────────────────

\section{Conclusão}

O desenvolvimento do LudiCode permitiu aplicar na prática conceitos estudados na disciplina de Compiladores. O projeto mostra como uma linguagem de programação pode ser criada a partir da definição de tokens, regras gramaticais e mecanismos de execução. Durante o trabalho, foram implementadas etapas fundamentais, como análise léxica, análise sintática, construção da AST e interpretação dos comandos. Cada uma dessas etapas possui uma função específica dentro do funcionamento do sistema.
Além do aspecto técnico, o projeto também possui uma proposta educacional, pois utiliza comandos em português para facilitar o aprendizado de programação. Dessa forma, o LudiCode demonstra como conceitos de compiladores podem ser aplicados em uma ferramenta voltada ao ensino de lógica computacional e robótica.

\end{document}

